\documentclass[11pt]{article}
\usepackage[margin=1in]{geometry}
\usepackage{amsmath,amssymb,amsthm}
\usepackage{xurl}
\usepackage{hyperref}
\sloppy
\let\oldus\_ \renewcommand{\_}{\oldus\allowbreak}
\newtheorem{theorem}{Theorem}
\newtheorem{lemma}[theorem]{Lemma}
\theoremstyle{definition}
\newtheorem{definition}[theorem]{Definition}
\title{Conjecture 00000002154 is false:\\ every cycle $C_m$ is closed Eulerian, but its critical group maps onto $\mathbb Z/m$}
\author{Jackmeson1}
\date{2026-10-05}
\begin{document}
\maketitle

\section{The conjecture}
Conjecture 00000002154 (English text, verbatim): ``Definition: The number of cyclic factors of the
critical group Jac(G) (the sandpile group) is the number of its cyclic Z-components. Conjecture:
The classification of graphs with zero cyclic factors (i.e. Jac entirely cyclic-free): exactly
subdivisions of closed Eulerian graphs (trivial exceptions exhausted).'' The Chinese text says the
same.

\paragraph{Reading.} Graphs are finite simple graphs. For a \emph{connected} graph, $\mathrm{Jac}(G)$ is a finite abelian group (for a disconnected graph $\mathrm{Div}^0(G)/\mathrm{im}\,L$ can be infinite, e.g.\ $\mathbb Z$ for two isolated vertices; every graph used below is connected);
its ``cyclic factors'' are the nontrivial cyclic summands in a decomposition of $\mathrm{Jac}(G)$
as a direct sum of cyclic groups. ``Zero cyclic factors'' (the text's gloss: ``Jac entirely
cyclic-free'') therefore means $\mathrm{Jac}(G)=0$. The conjecture asserts
\[ \mathrm{Jac}(G)=0 \iff G \text{ is a subdivision of a closed Eulerian graph}. \]
We refute the direction $\Leftarrow$ for the infinite family of all cycles $C_m$, $m\ge 3$. These
include $C_3$, the smallest closed Eulerian simple graph with an edge, so they are not ``trivial
exceptions''.

\paragraph{Not refuted.} If ``cyclic $\mathbb Z$-components'' meant summands isomorphic to the
infinite cyclic group $\mathbb Z$, then every connected graph would have zero of them, because
$\mathrm{Jac}(G)$ is finite. The conjecture would then say that every connected graph is a
subdivision of a closed Eulerian graph. That is also false, for example for $K_2$, but this
reading is not formalized here.

\section{Definitions}
For a finite vertex set $V$ let $L=D-A$ be the integer Laplacian of $G$, which is Mathlib's
\texttt{SimpleGraph.lapMatrix}.
\begin{definition}
$\mathrm{Div}^0(G)=\{x\in\mathbb Z^V:\sum_v x_v=0\}$, $\mathrm{Prin}(G)=L\,\mathbb Z^V$
(contained in $\mathrm{Div}^0$, Lean \texttt{prin\_le\_divZero}), and
$\mathrm{Jac}(G)=\mathrm{Div}^0(G)/\mathrm{Prin}(G)$.
For a sink $q\in V$ let $\Delta'$ be $L$ with the row and column of $q$ deleted. The second standard
model is the cokernel $\mathbb Z^{V\setminus q}/\Delta'\mathbb Z^{V\setminus q}$.
\end{definition}
Wikipedia, \emph{Abelian sandpile model} (\url{https://en.wikipedia.org/wiki/Abelian_sandpile_model}),
calls the second model ``the cokernel of the reduced graph Laplacian'' and says that it ``is most
commonly referred to as the sandpile group. Other common names for the same group are critical
group, Jacobian group or (less often) Picard group.'' We prove the result for both models.

\begin{definition}
A walk is an \emph{Eulerian circuit} if it is closed and uses every edge exactly once. Wikipedia,
\emph{Eulerian path} (\url{https://en.wikipedia.org/wiki/Eulerian_path}): ``an Eulerian circuit or
Eulerian cycle is an Eulerian trail that starts and ends on the same vertex.'' A graph is
\emph{closed Eulerian} if it is connected and has an Eulerian circuit. Wikipedia mentions a second
meaning of ``Eulerian graph'': ``a graph with every vertex of even degree''. $C_m$ satisfies this
one too, since every vertex has degree $2$.
\end{definition}

\begin{definition}
Wikipedia, \emph{Homeomorphism (graph theory)}: ``The subdivision of some edge $e$ with endpoints
$\{u,v\}$ yields a graph containing one new vertex $w$, and with an edge set replacing $e$ by two
new edges, $\{u,w\}$ and $\{w,v\}$.'' $H$ is a \emph{subdivision} of $G$ if $H$ is isomorphic to a
graph obtained from $G$ by finitely many such steps, possibly none.
\end{definition}

\section{The counterexample}
Let $C_m$ ($m\ge 3$) be the cycle on $\mathbb Z/m$, with $i\sim i\pm1$ (Mathlib's \texttt{cycleGraph m}).

\begin{lemma}\label{lem:euler}
$C_m$ is closed Eulerian, hence a subdivision (with no subdivided edge) of a closed Eulerian graph.
\end{lemma}
\begin{proof}
$C_m$ is connected. Mathlib's closed walk \texttt{cycleGraph.cycle} is a cycle of length $m$, so it
is a trail with $m$ distinct edges. All $m$ degrees equal $2$, so by the handshake lemma $C_m$ has
exactly $m$ edges, and the walk uses each of them exactly once.
\end{proof}

\begin{lemma}\label{lem:phi}
Define $\varphi:\mathbb Z^{\mathbb Z/m}\to\mathbb Z/m$ by $\varphi(x)=\sum_i i\,x_i \bmod m$. Then
$\varphi(Lf)=0$ for every $f$, and $\varphi(\delta_1-\delta_0)=1$.
\end{lemma}
\begin{proof}
$(Lf)_i=2f_i-f_{i-1}-f_{i+1}$. Reindexing, $\varphi(Lf)=\sum_j f_j\,(2j-(j+1)-(j-1))=0$ in
$\mathbb Z/m$. The second claim is immediate.
\end{proof}

\begin{theorem}\label{thm:main}
For every $m\ge3$, $\varphi$ induces a surjection $\mathrm{Jac}(C_m)\to\mathbb Z/m$. The
restriction $\varphi'(y)=\sum_{i\ne0} i\,y_i$ induces a surjection
$\mathbb Z^{V\setminus 0}/\Delta'\mathbb Z^{V\setminus 0}\to\mathbb Z/m$. In particular
$\mathrm{Jac}(C_m)\neq0$, and every decomposition of $\mathrm{Jac}(C_m)$ as a direct sum has a
nontrivial summand. Together with Lemma~\ref{lem:euler}, the class ``subdivisions of closed
Eulerian graphs'' contains graphs with a nonzero number of cyclic factors, so the classification is
false.
\end{theorem}
\begin{proof}
By Lemma~\ref{lem:phi}, $\varphi$ kills $\mathrm{Prin}$, and the class of $\delta_1-\delta_0$ maps to
the generator $1$. For the reduced model, extend $f\in\mathbb Z^{V\setminus0}$ by $0$ at the sink.
Then $(\Delta'f)_i=(L\tilde f)_i$ for $i\neq0$, and the weight of vertex $0$ is $0$, so
$\varphi'(\Delta'f)=\varphi(L\tilde f)=0$; the class of $\delta_1$ maps to $1$. If all summands of
a direct sum were trivial, the sum would be trivial, which contradicts $\mathrm{Jac}(C_m)\ne0$.
\end{proof}
\noindent\emph{Remark (not formalized).} By the same Wikipedia article, ``the order of the sandpile
group is the determinant of $\Delta'$, which by the matrix tree theorem is the number of spanning
trees''. $C_m$ has $m$ spanning trees, so the surjection is an isomorphism:
$\mathrm{Jac}(C_m)\cong\mathbb Z/m$.

\section{Lean formalization}
File \texttt{lean/Conjecture2154/Basic.lean} (324 lines, only \texttt{import Mathlib}), namespace
\texttt{C2154}.
\begin{itemize}
\item \texttt{divZero}, \texttt{prin} (range of \texttt{mulVecLin (lapMatrix Z G)}), \texttt{prin\_le\_divZero},
\texttt{Jac G := divZero V / (prin G).comap (divZero V).subtype}.
\item \texttt{redLap G q} (submatrix of \texttt{lapMatrix}), \texttt{SandpileGroup G q} (its cokernel).
\item \texttt{subdivideEdge H u v} (new vertex \texttt{none : Option W}), the inductive
\texttt{IsSubdivision} (constructors \texttt{refl}, \texttt{step}, \texttt{iso}), \texttt{IsClosedEulerian}
(\texttt{Connected} and a closed \texttt{Walk.IsEulerian}), \texttt{IsSubdivOfClosedEulerian}.
\item \texttt{cycleGraph\_isClosedEulerian}, \texttt{cycleGraph\_isSubdivOfClosedEulerian} (Lemma~\ref{lem:euler}).
\item \texttt{lap\_cycle\_apply}, \texttt{phi\_lap}, \texttt{phi\_witness} (Lemma~\ref{lem:phi});
\texttt{jacToZMod\_surjective}, \texttt{sandpileToZMod\_surjective}, \texttt{jac\_cycle\_nontrivial},
\texttt{jac\_cycle\_has\_cyclic\_factor} (Theorem~\ref{thm:main}).
\item Main: \texttt{not\_classification : forall n, IsSubdivOfClosedEulerian (cycleGraph (n+3)) /\textbackslash{}
Nontrivial (Jac (cycleGraph (n+3))) /\textbackslash{} (surjection Jac -> ZMod (n+3)) /\textbackslash{}
(surjection SandpileGroup -> ZMod (n+3))}; and \texttt{not\_conjecture}.
\end{itemize}
Axiom audit: \texttt{\#print axioms} for \texttt{not\_classification}, \texttt{not\_conjecture} and
\texttt{jac\_cycle\_has\_cyclic\_factor} each reports \texttt{[propext, Classical.choice, Quot.sound]}.
No \texttt{sorry}, no \texttt{native\_decide}, no added axioms.

\section{Relation to the earlier submission}
PR \#267 (by orionsheep) used $C_6$, a subdivision of $C_3$, with $\mathrm{Jac}(C_6)\cong\mathbb Z/6$.
Its Lean theorem proved only
$3\cdot2=6\wedge 2\bmod 2=0\wedge 6=2\cdot3\wedge 6>1\wedge2\ne0\wedge3\ne0$. The review rejected
it because the decisive content ``(Jac(C6) = Z/6 via Matrix--Tree and Smith normal form, and that
C6 subdivides the Eulerian C3) is prose- and script-only, with `6 > 1' standing in for
nontriviality by fiat.'' This submission defines the critical group in Lean (both
$\mathrm{Div}^0/\mathrm{Prin}$ and the reduced-Laplacian cokernel, from Mathlib's Laplacian), the
subdivision relation, and closed Eulerian graphs. It proves nontriviality through an explicit
surjection onto $\mathbb Z/m$ for every cycle $C_m$, $m\ge3$, not for a single graph. It uses the
trivial subdivision ($C_m$ itself), which is already in the class; the isomorphism
$\mathrm{Jac}(C_m)\cong\mathbb Z/m$ is not needed and is left as a remark.
\end{document}
